% Subespacios, producto cuña y determinante (partición por clase: apunte/scripts/reorganizar_por_clase.py)
\chapter{Subespacios, producto cuña y determinante}\label{cap:subespacios}\label{sec:bloque-algebraico}

\section{La idea}


Primero el mapa: cada operación del álgebra se lee como una pregunta geométrica. Lo armamos así y dejamos las cuentas para la sección siguiente.

\subsection{Del álgebra a la geometría: qué pregunta cada operación}\label{sec:algebra-geometria}

\textbf{Desarrollo.} El hilo de la materia es siempre el mismo: una operación algebraica,
en apariencia abstracta, resulta ser exactamente la pregunta geométrica que necesitábamos.
Conviene tener el mapa antes de las cuentas.

\begin{center}
\begin{tabularx}{\textwidth}{>{\bfseries}p{34mm} X X}
\toprule
Operación & Qué calcula & Qué pregunta geométrica responde\\
\midrule
Producto interno \(u\cdot v\) & un número & ¿son ortogonales? ¿cuánto se proyecta uno
sobre otro? En homogéneas: ¿el punto está en la recta?\\
Producto vectorial \(u\times v\) & un vector normal a ambos & ¿cuál es la recta que une dos
puntos? ¿cuál es el punto donde se cortan dos rectas? ¿qué área orientada forman?\\
Determinante \(\det[u\ v\ w]\) & un número & ¿están alineados tres puntos? ¿son
linealmente dependientes tres vectores? ¿qué volumen orientado forman?\\
Producto exterior \(u\wedge v\) & un bivector (área orientada) & lo mismo que el
vectorial, pero en cualquier dimensión\\
Núcleo de \(A\) & un subespacio & ¿qué vectores aplasta \(A\) a cero? Si hay alguno,
\(\det A=0\)\\
Autovectores de \(A\) & vectores con \(Av=\lambda v\) & ¿qué direcciones deja quietas
\(A\)? En homogéneas: ¿qué puntos deja fijos la transformación?\\
\bottomrule
\end{tabularx}
\end{center}

\begin{idea}[La cadena]
incidencia \(\to\) unión de puntos \(\to\) alineación \(\to\) determinante \(\to\)
núcleo \(\to\) autovectores \(\to\) puntos fijos. Cada eslabón se deduce del anterior;
es lo que desarrollan las secciones siguientes de este capítulo.
\end{idea}

\begin{idea}[El producto cuña como detector de colinealidad]
El producto cuña, detector de colinealidad, sirve primero para construir un polinomio cuyas
raíces son los autovalores, y luego para hallar los autovectores como el subespacio nulo
de unas matrices que armamos.
\end{idea}

\section{Construcciones, cuentas y ejemplos}


Hacemos la cuenta de cada operación con su figura: producto interno, producto vectorial, determinante, núcleo, autovalores y formas lineales. El orden es el de la cadena de la sección anterior.

El objetivo es que la parte geométrica y la parte algebraica se conecten, y entender por qué
el determinante tiene las propiedades ``mágicas'' que tiene. Lo que sigue reconstruye ese
razonamiento en orden.

\subsection{Tres puntos alineados: el determinante}\label{sec:determinante}

{\small\vspace{1mm}\linkUnidad{07a_cuna_determinante}{Simulación interactiva: Unidad 9 — Cuña y determinante: el detector de colinealidad}\par}


\textbf{Desarrollo.} ¿Está el punto \(r\) en la recta que une \(p\) y \(q\)? Se forma la
recta y se pregunta por la incidencia:
\[
  (p\times q)\cdot r=0 .
\]
Su anulación es simétrica en los tres puntos (si \(r\) está en \(pq\), entonces \(p\) está
en \(qr\) y \(q\) en \(rp\)); el valor, en cambio, es antisimétrico: se anula si dos de
ellos coinciden, es lineal en cada uno y cambia de signo al intercambiar dos. Es, término a término, el determinante:
\[
  (p\times q)\cdot r=\det\begin{pmatrix}p_1&p_2&p_3\\q_1&q_2&q_3\\r_1&r_2&r_3\end{pmatrix}.
\]

\begin{resultado}[El determinante es la única operación posible]
Una operación de tres vectores de \(\R^3\) que sea lineal en cada argumento y que se anule
cuando dos argumentos coinciden es, a menos de una constante, el determinante. Vale cero
exactamente cuando los tres vectores son linealmente dependientes (coplanares como vectores;
alineados como puntos homogéneos) y su valor absoluto es el volumen del paralelepípedo que
forman. Intercambiar dos filas cambia el signo: es la antisimetría otra vez.
\end{resultado}

\paragraph{Verificación numérica.} Con \(p=(1,2,3)\), \(q=(4,5,-1)\),
\(r=(0,0,1)\):
\[
  p\times q=\begin{pmatrix}2\cdot(-1)-3\cdot5\\ 3\cdot4-1\cdot(-1)\\ 1\cdot5-2\cdot4\end{pmatrix}
  =\begin{pmatrix}-17\\13\\-3\end{pmatrix},\qquad
  (p\times q)\cdot r=-3,
\]
y el determinante de la matriz de filas \(p,q,r\), desarrollado por la tercera fila, es
\(1\cdot(1\cdot5-2\cdot4)=-3\). Coinciden, como debía ser: desarrollar por una fila
\emph{es} hacer un producto vectorial con las otras dos y luego un producto interno.

\subsection{Núcleo de una matriz y determinante nulo}\label{sec:nucleo}

{\small\vspace{1mm}\linkUnidad{07b_nucleo_autovectores}{Simulación interactiva: Unidad 10 — Núcleo y autovectores: las direcciones que no giran}\par}


\textbf{Desarrollo.} El \emph{núcleo} de \(A\) es el conjunto de vectores \(v\) con
\(Av=0\). Si se piensa \(Av\) fila por fila, cada coordenada del resultado es el producto
interno de una fila de \(A\) con \(v\). Que las tres den cero significa que \(v\) es
ortogonal a las tres filas: las tres filas están en el plano de normal \(v\), son linealmente
dependientes, y por lo tanto
\[
  \boxed{\ \text{núcleo}(A)\neq\{0\}\quad\Longleftrightarrow\quad\det A=0\ }.
\]
Es la propiedad del determinante que se usa para encontrar autovalores.

\section{En profundidad}


Extendemos lo anterior con el producto exterior en cualquier dimensión y con el cambio de base y la diagonalización.

\subsection{Producto exterior: la misma idea en cualquier dimensión}\label{sec:exterior}

\textbf{Desarrollo.} El producto vectorial existe sólo en \(\R^3\). La operación que
generaliza la idea a cualquier dimensión es el \emph{producto exterior} de Grassmann,
\(u\wedge v\), definido únicamente por linealidad y antisimetría (\(u\wedge u=0\)). Su
resultado no es un vector sino un \emph{bivector}: un elemento de área orientada.

Escribiendo \(u=ae_1+be_2+ce_3\) y \(v=de_1+ee_2+fe_3\) y distribuyendo, los términos
\(e_i\wedge e_i\) se anulan y \(e_2\wedge e_1=-e_1\wedge e_2\), de modo que
\[
  u\wedge v=(ae-bd)\,e_1\wedge e_2+(af-cd)\,e_1\wedge e_3+(bf-ce)\,e_2\wedge e_3 .
\]
Son los mismos tres números del producto vectorial; sólo cambia el signo del término del
medio, porque la regla de la mano derecha asigna \(e_1\times e_3=-e_2\). Los ``planos
elementales'' \(e_1\wedge e_2\), \(e_1\wedge e_3\), \(e_2\wedge e_3\) son la base de las
áreas orientadas de \(\R^3\).

Operando con un tercer vector \(w=ge_1+he_2+ie_3\) sólo sobreviven los términos que completan
las tres direcciones, y reordenando (cada intercambio de vecinos cambia el signo) todo cae
sobre un único elemento de volumen:
\[
  (u\wedge v)\wedge w=\det[u\ v\ w]\;e_1\wedge e_2\wedge e_3 .
\]
De ahí la regla de los signos \(+,-,+\) del desarrollo por menores.

\begin{idea}[Contar subespacios]
En \(\R^3\) hay \(3\) direcciones de vectores, \(3\) planos elementales y \(1\) volumen; por
eso el resultado final es un número. En \(\R^4\) hay \(4\) direcciones, \(6\) planos, \(4\)
volúmenes tridimensionales y \(1\) hipervolumen. La operación es mecánica en cualquier
dimensión: da una fórmula cerrada para un vector ortogonal a un hiperplano de \(\R^n\), útil
para autovectores en dimensión baja cuando \(\operatorname{rango}(A-\lambda I)=n-1\). Para
matrices grandes los autovectores se calculan con otros algoritmos (QR, Lanczos).
\end{idea}

\paragraph{Ortogonal es lo que maximiza el volumen.} Entre todos los vectores de norma
\(1\), el que forma mayor volumen con un plano dado es el perpendicular al plano; los otros
producen un prisma ``inclinado'' de menor altura. De esta idea sale
la regla para encontrar el vector ortogonal a un subespacio en cualquier dimensión.
